% Context from chapters/machine-learning.tex; for mathematical reference, not compilation.
\quad c_\ell \leftarrow \frac{
		\sum_{i : y_i=\ell} x_i
	}{
		|\enscond{i}{y_i=\ell}|
	}
}
If a cluster becomes empty, its centroid does not contribute to the objective for the current labels. One may leave it in place or reseed it at a data point before the next assignment step. The centroid formula applies only to nonempty clusters.

Each step decreases or preserves $\Ee$. With a fixed tie-breaking rule that retains existing assignments at ties, and without indefinite reseeding, a changed assignment strictly decreases the objective. There are finitely many label assignments, so the algorithm terminates after finitely many such changes.

Because the objective is nonconvex, $k$-means need not find globally optimal clusters.
 
A reseeding heuristic moves centroids $c_\ell$ from redundant or low-contribution clusters to regions with high distortion. This can help the iterations escape a poor local configuration.

Figure~\ref{fig-kmeans} shows an example of $k$-means iterations on the Iris dataset.

\begin{figure}
\centering
\begin{tabular}{@{}c@{\hspace{10mm}}c@{}}
\includegraphics[width=.4\linewidth]{ml/pca-nn/kmeans-iter}&
\includegraphics[width=.35\linewidth]{ml/pca-nn/kmeans-classif-scores}
\end{tabular}
\caption{\label{fig-kmeans}
Left: iterations of the $k$-means algorithm. Right: class histograms after the $k$-means optimization.
}
\end{figure}


    
\paragraph{\texorpdfstring{$k$}{k}-means++}

The result of $k$-means depends strongly on initialization. Well-separated initial centers tend to cover the data more effectively than tightly grouped centers.

The randomized $k$-means++ initialization has an approximation guarantee even before Lloyd iterations are applied. In practice, Lloyd iterations usually improve the resulting centers. First choose $c_1$ uniformly among the sample points. After choosing $c_1,\ldots,c_\ell$, select $c_{\ell+1}$ from the sample according to a probability $\pi^{(\ell)}$ on $\{1,\ldots,n\}$ proportional to the squared distance from the nearest selected center
\eq{
	\foralls i \in \{1,\ldots,n\}, \quad 
	\pi^{(\ell)}_i \eqdef \frac{d_i^2}{\sum_{j=1}^n d_j^2}
	\qwhereq
	d_i\eqdef\min_{1\leq r\leq\ell}\norm{x_i-c_r}.
}
Points far from the previously seeded centers are more likely to be chosen. If every $d_i=0$, the current centers already achieve zero distortion and the procedure can stop.

The following theorem, due to David Arthur and Sergei Vassilvitskii, bounds the expected seeding cost by a logarithmic factor times the optimal cost. Computing a global optimum is NP-hard in general.

\begin{thm}
	Let $c^\star$ be the centroids selected by $k$-means++ and let $y^\star$ be their nearest-centroid labels, as in~\eqref{eq-kmeans-1}. Then
	\eq{
		\EE(\Ee(y^\star,c^\star)) \leq 8 (2+\log(k)) \umin{(y,c)}\Ee(y,c),
	}
	where the expectation is over the random choices made during initialization.
\end{thm}




    
\paragraph{Lloyd algorithm and continuous densities.}

The $k$-means iterations are also known as Lloyd's algorithm and are used in vector quantization for compression. The same alternating procedure applies to a probability measure $\mu$ on $\RR^p$ with finite second moment.
 
The energy to minimize becomes
\eq{
	\umin{ (\Vv,c) } \sum_{\ell=1}^k \int_{\Vv_\ell} \norm{x-c_\ell}^2 \d \mu(x)
}
where $(\Vv_\ell)_\ell$ is a partition of the domain. 
 
Step~\eqref{eq-kmeans-1} is replaced by the computation of a Voronoi cell
\eq{
	\foralls \ell \in \{1,\ldots,k\}, \quad 
	\Vv_\ell \eqdef \enscond{x}{ \foralls \ell' \neq \ell, \norm{x-c_\ell} \leq \norm{x-c_{\ell'}} }.
}
For distinct centroids, the Voronoi cells are polyhedra bounded by perpendicular bisectors. Resolve distance ties consistently to obtain a partition; this matters if $\mu$ assigns positive mass to a boundary. In low dimensions, computational geometry provides efficient algorithms for constructing the cells.
 
Step~\eqref{eq-kmeans-2} is then replaced by
\eq{
	\foralls \ell \in \{1,\ldots,k\}, \quad c_\ell \leftarrow \frac{
		\int_{\Vv_\ell} x \d\mu(x)
	}{
		\int_{\Vv_\ell} \d\mu(x)
	}.
}
The centroid update applies to cells of positive $\mu$-mass; zero-mass cells require a separate convention, as in the discrete algorithm. For a uniform planar density, hexagonal Voronoi patterns are characteristic of high-resolution quantization away from the boundary. They need not be exact finite-$k$ minimizers on a bounded domain.
 
Figure~\ref{fig-lloyd} displays two examples of Lloyd iterations on 2-D densities on a square domain.

\begin{figure}
\centering
\includegraphics[width=\linewidth]{ml/pca-nn/lloyd}
\caption{\label{fig-lloyd}
Continuous $k$-means (Lloyd) iterations for densities $\mu$. Top: uniform density. Bottom: nonuniform density, shown in grayscale as the density of $\mu$ with respect to Lebesgue measure.
}
\end{figure}




                                                  
                                                  
                                                  
\section{Empirical Risk Minimization}

We first describe empirical risk minimization, a framework shared by regression and classification. For classification, predictions may represent class probabilities rather than labels.

To make fitting $y_i \approx f(x_i)$ computationally tractable and obtain useful predictions, we restrict the admissible functions to a class of controlled complexity. The class may become richer as the sample size $n$ grows, but its complexity must remain compatible with the available data.

                                                  
\subsection{Empirical Risk}

Let $\Ff_n$ be a class of functions that may depend on the sample size. A common learning procedure is empirical risk minimization (ERM)
\eql{\label{eq-erm-1}
	\hat f \in \uargmin{f \in \Ff_n} \frac{1}{n} \sum_{i=1}^n \loss(f(x_i),y_i).
}
The loss function $\loss : \Yy^2 \rightarrow \RR^+$ measures prediction error for the chosen task.  
Classification and regression typically require different loss functions.
 
We sometimes write $\hat f_n$ to emphasize the estimator's dependence on the sample size $n$.


                                                  
\subsection{Prediction and Consistency}

For the analysis, assume the training pairs $(x_i,y_i)$ are i.i.d. with distribution $\pi$ on $\Xx \times \Yy$. The corresponding population optimization problem is
\eql{\label{eq-consistency-estim}
	\bar f \in \uargmin{f \in \Ff_\infty} \int_{\Xx \times \Yy} \loss(f(x),y) \d\pi(x,y) = \EE_{(\xp,\yp)\sim\pi}\bigl(\loss(f(\xp),\yp)\bigr).
}
The desired consistency property is $\hat f_n \rightarrow \bar f$ as $n \rightarrow +\infty$. One way to express this is through expected prediction error over the training samples $(x_i,y_i)_i$:
\eq{
	E_n \eqdef \EE( \tilde\loss(\hat f_n(\xp),\bar f(\xp)) ) \longrightarrow 0.
}
Here the expecta